\documentclass[11pt]{article}
\usepackage[margin=1in]{geometry}
\usepackage{amsmath,amssymb,amsthm,mathtools}
\usepackage{graphicx}
\usepackage{microtype}
\usepackage[colorlinks=true,linkcolor=blue,citecolor=blue,urlcolor=blue]{hyperref}

\newtheorem{theorem}{Theorem}
\newtheorem{lemma}{Lemma}
\newcommand{\E}{\mathbb E}
\newcommand{\Pp}{\mathbb P}
\newcommand{\R}{\mathbb R}
\newcommand{\one}{\mathbf 1}

\title{Finite-facet bodies and the centered-simplex test case for a Minkowski-sum gauge question}
\author{Candidate partial result for arXiv:1811.12548}
\date{22 July 2026}

\begin{document}
\maketitle

\begin{center}
\fbox{\parbox{0.88\textwidth}{\textbf{Status.}
Candidate partial result, likely valid.  This packet proves the conjectured
order for centered polytopes with linearly many facets and determines the
exact first-order asymptotic for centered simplices.  It does not settle the
universal worst-case question over all centered convex bodies.}}
\end{center}

\section{Source question}

Huang, Slomka, Tkocz, and Vritsiou ask in Section 6 of
\cite{HSTV} whether, for independent uniform random vectors $X,Y$ on a
centered convex body $K\subset\R^n$, there is a universal $\alpha<1$ for
which
\[
  \E\|X+Y\|_K\le 2-\Omega(n^{-\alpha}),
\]
and especially whether the worst deficit is of order $n^{-1/2}$.  Here
$\|z\|_K=\inf\{t>0:z\in tK\}$ is the Minkowski functional; it need not be a
norm when $K$ is nonsymmetric.  The source computes for the cube $Q_n$ that
\begin{equation}\label{eq:cube}
 \E\|X+Y\|_{Q_n}
 =2-\frac{2\,4^n}{(2n+1)\binom{2n}{n}}
 =2-\sqrt{\frac{\pi}{n}}+o(n^{-1/2}).
\end{equation}

Figure~\ref{fig:question} reproduces the complete source question from page
15.

\begin{figure}[ht]
  \centering
  \includegraphics[width=\textwidth]{figures/open_problem_crop.png}
  \caption{The question and cube benchmark in Section 6 of \cite{HSTV}.}
  \label{fig:question}
\end{figure}

\section{Result}

\begin{theorem}[Finite-facet bound]\label{thm:facets}
Let $K\subset\R^n$ be a polytope whose barycenter is the origin and which has
$N$ facets.  If $X,Y$ are independent and uniform on $K$, then
\begin{equation}\label{eq:facets}
 \E\|X+Y\|_K\le 2-\frac{3}{16\sqrt N}.
\end{equation}
Consequently, for every fixed $A>0$, the source's proposed
$n^{-1/2}$ deficit holds uniformly over centered polytopes with at most
$An$ facets.
\end{theorem}

\begin{theorem}[Centered-simplex asymptotic]\label{thm:main}
Let $K_n\subset\R^n$ be any simplex whose barycenter is the origin, and let
$X,Y$ be independent and uniform on $K_n$.  Then
\begin{equation}\label{eq:main}
 \E\|X+Y\|_{K_n}
 =2-\sqrt{\frac{\pi}{2n}}+o(n^{-1/2}).
\end{equation}
Consequently, the ratio of the simplex deficit to the cube deficit tends to
$1/\sqrt2$.  In particular, for all sufficiently large $n$,
\[
 \E\|X+Y\|_{K_n}>\E\|X+Y\|_{Q_n}.
\]
\end{theorem}

\subsection*{Proof intuition}

For a polytope the gauge is the maximum of its normalized facet
functionals.  The deficit from a fixed facet is a nonnegative log-concave
random variable of mean one, and such a variable has probability at most
$2t$ of lying in $[0,t]$.  For the sum of two independent points to have
gauge within $t$ of two, both points must lie within height $t$ of a common
facet.  The resulting probability is at most $4Nt^2$, yielding
Theorem~\ref{thm:facets} at $t\asymp N^{-1/2}$.

Affine invariance lets us put the simplex in barycentric coordinates.  A
uniform point is then a flat Dirichlet vector shifted so that the barycenter
is zero, and the simplex gauge is simply minus the smallest coordinate.
The problem therefore becomes an extreme-value calculation for
$\min_i(P_i+Q_i)$.

Flat Dirichlet coordinates are normalized independent exponentials.  Their
normalizing sums concentrate, so the relevant minimum is asymptotically the
minimum of $n+1$ independent $\mathrm{Gamma}(2,1)$ variables.  The
$\mathrm{Gamma}(2,1)$ distribution has quadratic lower tail $z^2/2$;
therefore its minimum has Rayleigh scaling and limiting mean
$\sqrt{\pi/2}$.  A uniform tail bound and an exponentially small
normalization bad event justify convergence of expectations.  The two
theorems still supply no universal comparison for bodies with arbitrary or
infinitely many supporting facets.

\section{Proof}

\subsection{The finite-facet theorem}

We first isolate the one-dimensional estimate.

\begin{lemma}[Log-concave boundary small ball]\label{lem:smallball}
Let $U$ be a nonnegative random variable with a log-concave density and
$\E U=1$.  Then, for $0\le t\le1/2$,
\[
 \Pp(U\le t)\le2t.
\]
\end{lemma}

\begin{proof}
Let $S(t)=\Pp(U\ge t)$.  The survival function of a log-concave density is
log-concave, so $H=-\log S$ is convex and $H(0)=0$.  For $u\ge t>0$,
convexity gives $H(u)\ge(u/t)H(t)$.  If $H(t)=0$ the conclusion is
immediate, so assume $H(t)>0$.  Then
\[
 1=\E U=\int_0^\infty S(u)\,du
 \le t+\int_t^\infty e^{-(u/t)H(t)}\,du
 \le t+\frac{t}{H(t)}.
\]
Thus $H(t)\le t/(1-t)\le2t$, and consequently
\[
 \Pp(U\le t)=1-e^{-H(t)}\le H(t)\le2t.
\]
The endpoint $t=0$ follows by continuity.
\end{proof}

\begin{proof}[Proof of Theorem~\ref{thm:facets}]
Choose normalized facet functionals $\varphi_1,\ldots,\varphi_N$ such that
\[
 K=\{x:\varphi_j(x)\le1\text{ for }1\le j\le N\}.
\]
Since the origin is in the interior,
$\|z\|_K=\max_j\varphi_j(z)$.  Put
\[
 U_j=1-\varphi_j(X),\qquad V_j=1-\varphi_j(Y).
\]
Each variable is nonnegative.  The one-dimensional marginal of uniform
measure on a convex body is log-concave, and the barycenter assumption gives
$\E\varphi_j(X)=0$; hence $\E U_j=1$.  Lemma~\ref{lem:smallball} applies.

The following deficit identity is exact:
\begin{equation}\label{eq:facetdeficit}
 2-\|X+Y\|_K=\min_{1\le j\le N}(U_j+V_j)=:D.
\end{equation}
For $0\le t\le1/2$, independence and a union bound give
\begin{align*}
 \Pp(D\le t)
 &\le\sum_{j=1}^N\Pp(U_j+V_j\le t)\\
 &\le\sum_{j=1}^N\Pp(U_j\le t)\Pp(V_j\le t)
 \le4Nt^2.
\end{align*}
Take $t_0=1/(4\sqrt N)$.  Then $\Pp(D>t_0)\ge3/4$, so
\[
 \E D\ge t_0\Pp(D>t_0)\ge\frac{3}{16\sqrt N}.
\]
Together with \eqref{eq:facetdeficit}, this proves \eqref{eq:facets}.
\end{proof}

\subsection{The exact simplex asymptotic}

Put $m=n+1$ and work in the hyperplane
\[
 H_m=\{x\in\R^m:\textstyle\sum_{i=1}^m x_i=0\}.
\]
Every centered $n$-simplex is a linear image of
\[
 K=\{x\in H_m:x_i\ge -1\text{ for every }i\}.
\]
Both uniform measure and the Minkowski functional commute with invertible
linear maps, so it suffices to treat this $K$.

Let $P=(P_1,\ldots,P_m)$ and $Q=(Q_1,\ldots,Q_m)$ be independent and uniform
on the probability simplex
\[
 \Delta_{m-1}=\{p_i\ge0:\textstyle\sum_i p_i=1\}.
\]
Then $X=mP-\one$ and $Y=mQ-\one$ are independent and uniform on $K$.  For
$z\in H_m$,
\[
 \|z\|_K=\inf\{t>0:z_i\ge-t\ \forall i\}=-\min_i z_i.
\]
It follows exactly that
\begin{equation}\label{eq:gaugeidentity}
 \|X+Y\|_K=2-mM_m,
 \qquad
 M_m:=\min_{1\le i\le m}(P_i+Q_i).
\end{equation}

We use the exponential representation.  Let $E_i,F_i$ be mutually
independent exponential random variables of mean one, and set
\[
 S_m=\sum_{i=1}^m E_i,\qquad T_m=\sum_{i=1}^m F_i.
\]
Then
\[
 (P_i)_{i=1}^m\stackrel d=(E_i/S_m)_{i=1}^m,
 \qquad
 (Q_i)_{i=1}^m\stackrel d=(F_i/T_m)_{i=1}^m.
\]

First ignore the two random denominators.  The variables
$Z_i:=E_i+F_i$ are independent $\mathrm{Gamma}(2,1)$ variables, with survival
function
\[
 q(z):=\Pp(Z_i>z)=e^{-z}(1+z).
\]
Define
\[
 W_m:=\sqrt m\min_{1\le i\le m}Z_i.
\]
For every fixed $t\ge0$,
\begin{align*}
 \Pp(W_m>t)
 &=q(t/\sqrt m)^m\\
 &=\left(1-\frac{t^2}{2m}+O(m^{-3/2})\right)^m
 \longrightarrow e^{-t^2/2}.
\end{align*}
Thus $W_m$ converges in distribution to the Rayleigh random variable $W$
with survival function $\Pp(W>t)=e^{-t^2/2}$.

We record the uniform tail control needed below.  Since
\[
 z-\log(1+z)\ge z^2/4\quad(0\le z\le1)
\]
and
\[
 z-\log(1+z)\ge(1-\log2)z\quad(z\ge1),
\]
we have, with $c_0=1-\log2>0$,
\begin{equation}\label{eq:tails}
 \Pp(W_m>t)\le
 \begin{cases}
   e^{-t^2/4},&0\le t\le\sqrt m,\\
   e^{-c_0t\sqrt m},&t\ge\sqrt m.
 \end{cases}
\end{equation}
These bounds give a common integrable majorant for the survival functions
and also uniformly bound $\E W_m^2$.  Hence, by the survival-integral
formula and dominated convergence,
\begin{equation}\label{eq:Wmean}
 \E W_m\longrightarrow \E W
 =\int_0^\infty e^{-t^2/2}\,dt
 =\sqrt{\frac\pi2}.
\end{equation}

It remains to restore the denominators.  Put $\delta_m=m^{-1/4}$ and let
\[
 G_m=\left\{\left|S_m/m-1\right|\le\delta_m,
                 \ \left|T_m/m-1\right|\le\delta_m\right\}.
\]
Standard exponential Chernoff bounds, obtained directly from
$\E e^{\theta E_i}=(1-\theta)^{-1}$, give
\begin{equation}\label{eq:bad}
 \Pp(G_m^c)\le4e^{-c\sqrt m}
\end{equation}
for a universal $c>0$ and all sufficiently large $m$.  On $G_m$, coordinate
by coordinate,
\[
 \frac{E_i+F_i}{m(1+\delta_m)}
 \le \frac{E_i}{S_m}+\frac{F_i}{T_m}
 \le \frac{E_i+F_i}{m(1-\delta_m)}.
\]
Consequently, if
\[
 A_m:=m^{3/2}M_m,
\]
then on $G_m$
\begin{equation}\label{eq:sandwich}
 \frac{W_m}{1+\delta_m}\le A_m\le
 \frac{W_m}{1-\delta_m}.
\end{equation}

Finally, $M_m\le 2/m$ because the minimum of the $m$ numbers
$P_i+Q_i$ is no larger than their average.  Thus $A_m\le2\sqrt m$.
Combining this deterministic bound with \eqref{eq:bad}, and combining the
uniform second-moment bound for $W_m$ with Cauchy--Schwarz, gives
\[
 \E[A_m\one_{G_m^c}]\longrightarrow0,
 \qquad
 \E[W_m\one_{G_m^c}]\longrightarrow0.
\]
The sandwich \eqref{eq:sandwich}, $\delta_m\to0$, and \eqref{eq:Wmean}
therefore imply
\begin{equation}\label{eq:Amean}
 \E A_m\longrightarrow\sqrt{\frac\pi2}.
\end{equation}
Using \eqref{eq:gaugeidentity} and $m=n+1$, we obtain
\[
 2-\E\|X+Y\|_K
 =m\E M_m
 =\frac{\E A_m}{\sqrt m}
 =\sqrt{\frac{\pi}{2n}}+o(n^{-1/2}),
\]
which proves \eqref{eq:main}.  Comparing with \eqref{eq:cube} proves the
last assertions. \qed

\section{Verification and scope}

The accompanying script numerically integrates the exact survival function
$q(t/\sqrt m)^m$ and independently samples the normalized Dirichlet model.
Those computations check the scaling and constants but are not used in the
proof.

The proof establishes the desired order for centered polytopes with
$O(n)$ facets and the exact leading term for simplices.  It is consistent
with, but does not prove, a universal deficit of order $n^{-1/2}$ for bodies
with arbitrary or infinitely many supporting facets.  It also shows that the cube is not a maximizer of
$\E\|X+Y\|_K$ in high dimension, should one consider that stronger variant.
The published source passage itself asks for the order of the worst deficit
and does not require cube extremality.

A bounded novelty check searched the local run indexes, the parsed arXiv
corpus, arXiv, and exact web phrases involving uniform simplex points,
Minkowski gauges of $X+Y$, and minima of sums of flat Dirichlet coordinates.
No prior exact computation was found.  This is not an exhaustive literature
claim.

\begin{thebibliography}{9}
\bibitem{HSTV}
H. Huang, B. A. Slomka, T. Tkocz, and B.-H. Vritsiou,
\emph{Improved bounds for Hadwiger's covering problem via thin-shell
estimates}, arXiv:1811.12548, especially Section 6 (page 15).
\end{thebibliography}

\end{document}
